\documentclass[a4paper, 11pt]{article}
\usepackage[utf8]{inputenc}
\usepackage[left=20mm, right=20mm, top=22mm, bottom=22mm]{geometry}

\usepackage{amsmath,amsfonts,amssymb,bm}
\usepackage[english]{babel}
\usepackage{hyperref}
\usepackage{xcolor}
\usepackage{fancyhdr}
\usepackage{graphicx}
\usepackage{booktabs}
\usepackage{float}

\definecolor{harvardcrimson}{HTML}{8C1D40}

\hypersetup{
  colorlinks=true,
  linkcolor=harvardcrimson,
  citecolor=harvardcrimson,
  urlcolor=harvardcrimson,
  hypertexnames=false,
}

\pagestyle{fancy}
\fancyhf{}
\rhead{\textcolor{harvardcrimson}{\textbf{Optical projection on cluster lensing}}}
\lhead{\textcolor{harvardcrimson}{\textbf{J.\,H.\,Esteves}}}
\cfoot{\thepage}
\renewcommand{\headrulewidth}{0.4pt}
\renewcommand{\footrulewidth}{0.4pt}
\setlength{\headheight}{14pt}
\renewcommand{\baselinestretch}{1.15}
\setlength{\emergencystretch}{3em}
\sloppy

% Custom commands -----------------------------------------------------------
\newcommand{\lob}{\lambda^{\rm ob}}
\newcommand{\ltr}{\lambda^{\rm tr}}
\newcommand{\zob}{z^{\rm ob}}
\newcommand{\ztr}{z^{\rm tr}}
\newcommand{\dprj}{\Delta^{\rm prj}}
\newcommand{\xinl}{\xi_{\rm NL}}
\newcommand{\hMpc}{h^{-1}\,\mathrm{Mpc}}
\newcommand{\hMsun}{h^{-1}\,M_{\odot}}
\newcommand{\Pop}{\mathcal{P}}
\newcommand{\bsel}{b_{\rm sel}}
\newcommand{\bhalo}{b_{\rm halo}}
\newcommand{\bbhalo}{\avg{b_{\rm halo}}}
\newcommand{\bLOS}{b_{\rm sel}^{\rm LOS}}
\newcommand{\bLSS}{b_{\rm sel}^{\rm LSS}}
\newcommand{\bbsel}{\avg{b_{\rm sel}}}
\newcommand{\bbLOS}{\avg{b_{\rm sel}^{\rm LOS}}}
\newcommand{\bbLSS}{\avg{b_{\rm sel}^{\rm LSS}}}
\newcommand{\bbselij}{\avg{b_{{\rm sel},ij}}}
\newcommand{\bbLOSij}{\avg{b_{{\rm sel},ij}^{\rm LOS}}}
\newcommand{\bbLSSij}{\avg{b_{{\rm sel},ij}^{\rm LSS}}}
\newcommand{\bbhaloij}{\avg{b_{{\rm halo},ij}}}
\newcommand{\rhoprj}{\rho^{\rm prj}}
\newcommand{\Sprj}{\Sigma^{\rm prj}}
\newcommand{\Smis}{\Sigma_{\rm mis}}
\newcommand{\DSprj}{\Delta\Sigma^{\rm prj}}
\newcommand{\DSmis}{\Delta\Sigma_{\rm mis}}
\newcommand{\Rmis}{R_{\rm mis}}
\newcommand{\avg}[1]{\langle #1 \rangle}

\title{\vspace{-4ex}\huge \textbf{Effects of Optical Projection on
        Cluster Lensing}\\[0.5ex]
       \Large \textbf{a forward-model reference for the DES Y3
                      cluster-cosmology pipeline}}
\author{J.\,H.\,Esteves}
\date{2026-05-15}

\begin{document}
\maketitle

\begin{abstract}
Photometric cluster finders such as \textsc{redMaPPer} count richness
through red-sequence members along a single line of sight (LOS), and red
galaxies are biased tracers of large-scale structure: at fixed halo mass,
LOS aligned with filaments are assigned a larger observed richness $\lob$
than unstructured ones, and unrelated haloes inside the photo-$z$
window contaminate the member list.  Both effects bias the stacked
weak-lensing signal $\avg{\DSprj(R\,|\,\lob,\zob)}$ at the few-to-tens of
percent level and, if unmodelled, propagate directly into the inferred
mass--richness relation and cosmological parameters.  This document is
the reference paper for the projection-effects module of the DES Y3
cluster-cosmology pipeline: it consolidates the formalism of
Costanzi et~al.\ (2019, 2026), establishes a unified notation, derives
the forward model implemented in our code (the
\texttt{RichnessSelection} package), and lists the calibrated parameters
used by the production pipeline.  All radial integrals are kept in real
space; no Fourier-space short-cuts are used.
\end{abstract}

\tableofcontents
\vspace{1ex}\hrule\vspace{1ex}

% ============================================================================
\section{Introduction and scope}
\label{sec:intro}

The mean lensing surface density around optically selected clusters of
observed richness $\lob$ at observed redshift $\zob$ is the sum of a
one-halo term (the cluster's own miscentered NFW profile) and a
two-halo term (correlated neighbouring haloes along the LOS):
\begin{equation}
\avg{\Sigma(R\,|\,\lob,\zob)} \;=\;
\avg{\Sigma^{\rm 1h}(R\,|\,\lob,\zob)} \;+\;
\avg{\Sprj(R\,|\,\lob,\zob)} ,
\label{eq:sigma_split}
\end{equation}
where $\avg{\Sigma^{\rm 1h}}$ is the miscentered NFW of the target halo
and $\avg{\Sprj}$ is the two-halo projected surface density that the
\texttt{RichnessSelection} module computes.  The physical origin of the
selection effect is that, at fixed halo mass, a LOS with more red
galaxies receives a larger $\lob$ than one with fewer; the two-halo
term at fixed $\lob$ therefore does not carry the bare halo bias
$b(M,z)$ but a \emph{selection-affected} cluster bias
$\bbselij(\theta)$ which interpolates between
two physically distinct regimes: \emph{correlated structures in the
LOS} $\bbLOS$, set at small angular separations by close red-galaxy
neighbours projecting along the line of sight, and the \emph{LSS
regime} $\bbLSS$, set at large angular separations by correlated
large-scale structure.

This document is organised following the structure of
Park et~al.\ (2023): Sec.~\ref{sec:modelling} develops the full
forward model (halo-model ingredients, the projection-affected richness
kernel, the selection-affected cluster bias, the two-halo projected
lensing profile, miscentering, photo-$z$ scatter, and the parameter
inference setup); Sec.~\ref{sec:tests} validates the model against the
CAMB-based emulator and the numerical pipeline against
\texttt{scipy.quad}, and presents an end-to-end test on Buzzard mocks;
Sec.~\ref{sec:application} applies the model to the DES Y3
\textsc{redMaPPer} catalogue; Sec.~\ref{sec:discussion} contrasts our
real-space approach with the Park 2023 $\Pi(R)$ boost and the
Costanzi 2026 multiplicative profile fit; Sec.~\ref{sec:summary}
concludes.  All radial integrals are evaluated in real space, with no
Fourier-space short-cut and no global multiplicative fitting function
on the lensing profile.

% ============================================================================
\section{Modelling}
\label{sec:modelling}

\subsection{Cluster number counts and selection function}
\label{sec:counts}

The first observable of the DES Y3 cluster-cosmology pipeline is the
abundance of clusters in joint bins of observed richness and observed
redshift,
$\Delta\lambda_i\equiv[\lambda_i^{\min},\lambda_i^{\max}]$ and
$\Delta z_j\equiv[z_j^{\min},z_j^{\max}]$.  In the forward-model
formulation, the expected number of clusters in the $(i,j)$-th bin is
\begin{equation}
\avg{N_{ij}} \;=\;
\int dM \int d\ztr \;
\Omega(\ztr)\,\frac{dV}{d\Omega\,d\ztr}\,n(M,\ztr)\,
\mathcal{S}_{ij}(M,\ztr) ,
\label{eq:Nij}
\end{equation}
where $n(M,\ztr)=dn/dM$ is the halo mass function,
$dV/(d\Omega\,d\ztr)$ is the comoving volume element,
$\Omega(\ztr)$ is the effective survey solid angle, and
$\mathcal{S}_{ij}(M,\ztr)$ is the
\textsc{redMaPPer} \emph{selection function} --- the probability that
a halo of $(M,\ztr)$ is scattered into the observed bin.  Marginalising
over the latent richness $\ltr$ and the observed bins,
\begin{equation}
\mathcal{S}_{ij}(M,\ztr) \;=\;
\int_0^{\infty} d\ltr \;
P(\ltr\,|\,M,\ztr)
\int_{\Delta\lambda_i} d\lob
\int_{\Delta z_j} d\zob \;
P(\lob\,|\,\ltr,\ztr)\,
P(\zob\,|\,\ztr,\Delta\lambda_i) ,
\label{eq:Sel_full}
\end{equation}
where $P(\ltr\,|\,M,\ztr)$ is the mass--richness relation,
$P(\lob\,|\,\ltr,\ztr)$ is the projection-affected richness kernel of
Sec.~\ref{sec:richness}, and
$P(\zob\,|\,\ztr,\Delta\lambda_i)$ is the photo-$z$ kernel.  Because
$P(\zob\,|\,\ztr,\Delta\lambda_i)$ is well approximated by a Gaussian in
DES Y3, the $\zob$ integral factorises into a difference of normal
cdfs and the selection function separates into observed-richness and
observed-redshift pieces,
\begin{equation}
\mathcal{S}_{ij}(M,\ztr) \;=\;
\mathcal{S}_i(M,\ztr)\,
\mathcal{S}_j(\ztr) ,
\qquad
\mathcal{S}_j(\ztr) \;=\;
\Phi\!\left(\frac{\zob-\ztr}{\sigma_z}\right)\bigg|_{\Delta z_j} ,
\label{eq:Sel_split}
\end{equation}
with $\sigma_z\equiv\sigma_z(\Delta\lambda_i)$ the (richness-dependent)
photo-$z$ scatter and the evaluation bar in the usual
definite-integral sense.  The remaining \emph{richness selection
function} is
\begin{equation}
\mathcal{S}_i(M,\ztr) \;\equiv\;
\Pr(\lob\in\Delta\lambda_i\,|\,M,\ztr) \;=\;
\int_0^{\infty} d\ltr \;
P(\ltr\,|\,M,\ztr)\,
\mathcal{S}_{\Delta\lambda_i}(\ltr,\ztr) ,
\label{eq:Sel_lambda}
\end{equation}
the fraction of haloes at $(M,\ztr)$ that are scattered into
$\Delta\lambda_i$ by the observational kernel, with the bin-integrated
richness kernel
\begin{equation}
\mathcal{S}_{\Delta\lambda_i}(\ltr,\ztr) \;\equiv\;
\int_{\Delta\lambda_i} d\lob \;
P(\lob\,|\,\ltr,\ztr) .
\label{eq:Slam_def}
\end{equation}
Equation~\eqref{eq:Sel_lambda} is the form used inside
Eq.~\eqref{eq:Nij} and inside every binned lensing observable below.
The non-Gaussian projection tail of $P(\lob\,|\,\ltr,\ztr)$ derived
in Sec.~\ref{sec:richness} enters $\avg{N_{ij}}$ exclusively through
$\mathcal{S}_{\Delta\lambda_i}$, while the same kernel sets the
projection contamination on the lensing profile through the
selection-affected cluster bias of Sec.~\ref{sec:bias}.

% ----------------------------------------------------------------------------
\subsection{Halo-model predictions and notation}
\label{sec:halo_model}

Masses are in $\hMsun$; $R$ is a physical projected distance; $\theta$
is a sky angle; $\chi$, $R_\lambda$, $R_{\rm excl}$ are comoving in
$\hMpc$.  Default cosmology is Buzzard-v1.1
($\Omega_m=0.286$, $h=0.7$, $\sigma_8=0.82$).  Key shorthand:
\begin{equation}
R_\lambda(\lob)=(\lob/100)^{0.2}\,\hMpc,\qquad
R_{\rm excl}\equiv R_\lambda(\lob)(1+\zob),\qquad
\theta_\lambda \equiv R_\lambda(\lob)/D_A(\zob).
\label{eq:Rlam}
\end{equation}
The halo-model ingredients are the Tinker-2008 halo mass function
$n(M,z)$, the Tinker-2010 peak-height halo bias $b(M,z)$, and the
nonlinear matter correlation function $\xinl(r,z)$ obtained from a
halofit non-linear $P_{\rm NL}(k,z)$ via an FFTlog Hankel transform of
the linear $P(k,z)$ from CAMB.  Halo exclusion sets $\xinl\equiv 0$
whenever the 3-D comoving separation $\Delta\chi<R_{\rm excl}$.  The
mass--richness relation $P(\ltr\,|\,M,z)$ is the standard log-normal
of Costanzi et~al.\ (2019).  Photo-$z$ on the target, target
miscentering, triaxiality, and \textsc{redMaPPer} percolation
member-loss are deferred to Sec.~\ref{sec:miscentering}--%
\ref{sec:photoz} or to future work.

% ----------------------------------------------------------------------------
\subsection{Projection-affected richness kernel}
\label{sec:richness}

The forward model uses the Costanzi et~al.\ (2019) kernel for the
mapping from intrinsic to observed richness.  It starts from the
additive decomposition
\begin{equation}
\lob \;=\; \ltr + \Delta^{\rm bkg} + \dprj ,
\qquad
\Delta^{\rm bkg}\sim \mathcal{N}(\Delta\mu,\sigma^{2}) ,
\label{eq:lob_decomp}
\end{equation}
where $\Delta^{\rm bkg}$ is a Gaussian background/measurement
fluctuation and $\dprj\ge 0$ is a one-sided projection boost modelled
as a spike at zero plus an exponential tail:
\begin{equation}
P(\dprj\,|\,\ltr,z) \;=\;
(1-f^{\rm prj})\,\delta_D(\dprj)
\;+\;
f^{\rm prj}\,\tau\,e^{-\tau\dprj}\,\Theta(\dprj) .
\label{eq:prj_spike_tail}
\end{equation}
Convolving the Gaussian $\Delta^{\rm bkg}$ with the projection boost
$\dprj$ produces an \emph{exponentially modified Gaussian} (EMG)
distribution: the law of $X=G+E$ with $G\sim\mathcal{N}$ and
$E\sim\mathrm{Exp}$.  The full kernel is therefore
\begin{equation}
\begin{aligned}
P(\lob\,|\,\ltr,z) \;=\;&
(1-f^{\rm prj})\,\mathcal{N}(\lob;\,\mu,\sigma) \\
&{}+ f^{\rm prj}\,\frac{\tau}{2}\,
\exp\!\left[\frac{\tau}{2}\left(2\mu+\tau\sigma^{2}-2\lob\right)\right]
\operatorname{erfc}\!\left(\frac{\mu+\tau\sigma^{2}-\lob}{\sqrt{2}\,\sigma}\right) ,
\end{aligned}
\label{eq:emg_kernel}
\end{equation}
with $\mu\equiv\ltr+\Delta\mu$.  The four parameters
$\{\Delta\mu,\sigma,f^{\rm prj},\tau\}$ depend on $(\ltr,z)$ and are
calibrated empirically (synthetic-cluster injections in SDSS).  The
first line of Eq.~\eqref{eq:emg_kernel} is the background-only limit
recovered when $f^{\rm prj}=0$; the second line is the EMG projection
contribution.  The bare delta-plus-exponential form often quoted in the
literature (e.g.\ Costanzi et~al.\ 2026) omits the Gaussian background
convolution and is not used here.

The MOR pdf $P(\ltr\,|\,M,z)$ is taken to be the standard log-normal
of Costanzi et~al.\ (2019); together with Eq.~\eqref{eq:emg_kernel} it
defines $P(\lob\,|\,M,z)$ via
\begin{equation}
P(\lob\,|\,M,z) \;=\;
\int d\ltr \, P(\lob\,|\,\ltr,z)\,P(\ltr\,|\,M,z) ,
\label{eq:Plob_M}
\end{equation}
which is the only place where the projection tail enters the rest of
the pipeline.

\subsubsection{Closed form of the richness selection function}
\label{sec:richness:closedform}

Inserting the EMG kernel of Eq.~\eqref{eq:emg_kernel} into the
bin-integrated richness kernel $\mathcal{S}_{\Delta\lambda_i}(\ltr,z)$
of Eq.~\eqref{eq:Slam_def} and exchanging the bin integral with the
EMG decomposition gives
\begin{equation}
\mathcal{S}_{\Delta\lambda_i}(\ltr,z) \;=\;
(1-f^{\rm prj})\,\mathcal{S}_{\Delta\lambda_i}^{\rm G}(\ltr,z)
\;+\;
f^{\rm prj}\,\mathcal{S}_{\Delta\lambda_i}^{\rm EMG}(\ltr,z) ,
\label{eq:Slam_split}
\end{equation}
both pieces of which admit a closed form.  The Gaussian piece is the
ordinary bin-cdf difference,
\begin{equation}
\mathcal{S}_{\Delta\lambda_i}^{\rm G}(\ltr,z) \;=\;
\Phi\!\left(\frac{\lambda_i^{\max}-\mu}{\sigma}\right)
- \Phi\!\left(\frac{\lambda_i^{\min}-\mu}{\sigma}\right) ,
\qquad \mu \equiv \ltr+\Delta\mu ,
\label{eq:SlamG}
\end{equation}
where $\Phi$ is the standard normal cdf.  The EMG piece is the
difference of the EMG cdf evaluated at the bin edges,
\begin{equation}
\mathcal{S}_{\Delta\lambda_i}^{\rm EMG}(\ltr,z) \;=\;
F_{\rm EMG}(\lambda_i^{\max};\mu,\sigma,\tau)
- F_{\rm EMG}(\lambda_i^{\min};\mu,\sigma,\tau) ,
\label{eq:SlamEMG}
\end{equation}
with the closed-form EMG cdf
\begin{equation}
F_{\rm EMG}(x;\mu,\sigma,\tau) \;=\;
\Phi\!\left(\frac{x-\mu}{\sigma}\right)
- \exp\!\left[-\tau(x-\mu)+\tfrac{1}{2}\tau^{2}\sigma^{2}\right]
  \Phi\!\left(\frac{x-\mu}{\sigma}-\tau\sigma\right) .
\label{eq:exg_cdf}
\end{equation}
The first term is the Gaussian cdf at $x$; the second encodes the
correction introduced by the exponential projection tail.
Equation~\eqref{eq:exg_cdf} follows from convolving the Gaussian and
exponential laws and integrating against the bin: the
$\lob$-integration in Eq.~\eqref{eq:Sel_full} is therefore performed
\emph{analytically}, so that the richness selection function reduces
to the 1-D integral of Eq.~\eqref{eq:Sel_lambda} over the latent
variable $\ltr$ alone, weighted by the mass--richness pdf
$P(\ltr\,|\,M,z)$.  This closed form is the form actually evaluated in
the production code; the full derivation is given in
\texttt{docs/richness\_selection\_function.tex}.

% ----------------------------------------------------------------------------
\subsection{Selection-affected cluster bias}
\label{sec:bias}

\subsubsection{Halo-bias aggregate \texorpdfstring{$\bbhaloij$}{bhalo\_ij}}

The halo-bias aggregate is the bin-averaged halo bias weighted by the
halo-mass distribution and the selection function of
Eq.~\eqref{eq:Sel_lambda}:
\begin{equation}
\bbhaloij
 \;=\; \frac{\int dM \int d\ztr\;n(M,\ztr)\,b(M,\ztr)\,
             \mathcal{S}_{ij}(M,\ztr)}
            {\int dM \int d\ztr\;n(M,\ztr)\,
             \mathcal{S}_{ij}(M,\ztr)} .
\label{eq:bhalo}
\end{equation}
$n(M,\ztr)$ is the Tinker-2008 halo mass function, $b(M,\ztr)$ the
Tinker-2010 peak-height halo bias, and $\mathcal{S}$ is the
selection function of Eq.~\eqref{eq:Sel_full}.  By construction
$\bbhaloij$ is a function of the observed bin
$(\Delta\lambda_i,\Delta z_j)$ only.

\subsubsection{Mixture-of-regimes ansatz}

Optically selected clusters carry two physically distinct
contributions to the cluster--matter correlation: \emph{correlated
structures in the LOS}, dominant at small angular separations and
sourced by close red-galaxy neighbours projecting along the line of
sight into the target's \textsc{redMaPPer} aperture, and the
\emph{LSS regime}, set at large angular separations by correlated
large-scale structure.  We therefore write the
selection-affected cluster bias as a \emph{mixture model} over the
two regimes,
\begin{equation}
\bbselij(\theta) \;=\;
\bbLOSij\,[\,1-\sigma(\theta)\,]
\;+\;
\bbLSSij\,\sigma(\theta) ,
\qquad
\sigma(\theta) \;=\; \frac{1}{1+e^{-k(\theta-\theta_0)}} ,
\label{eq:b_sigmoid}
\end{equation}
with shape constants fixed analytically to $k=2.5/\theta_\lambda$ and
$\theta_0=\theta_\lambda/2$ (Costanzi et~al.\ 2026).  Throughout this
note, $\avg{\,\cdot\,}$ denotes averaging over the latent richness
$\ltr$ against the bin posterior of Eq.~\eqref{eq:plt_post}, so the
two amplitudes
\begin{equation}
\bbLOSij ,\qquad
\bbLSSij
\label{eq:b_regimes}
\end{equation}
are explicitly latent-marginalised quantities and depend only on the
\emph{observed} bin $(\Delta\lambda_i,\Delta z_j)$.  These are the
physically meaningful objects of the model.  The intermediate
$\ltr$-dependent quantities introduced in the pipeline below
(Eqs.~\eqref{eq:b_LSS_latent}--\eqref{eq:b_LOS_latent}) carry no
particular name and serve only as the integrand of the
$\ltr$-marginalisation.

\subsubsection{Projection kernel and \texorpdfstring{$\Pop[X]$}{P[X]} operator}

The \emph{projection kernel} $\rhoprj$ is the per-halo weight with
which a projected halo at $(\lambda,z,\theta)$ contaminates the
observed richness of a target at $(\lob,\zob)$,
\begin{equation}
\rhoprj(\lambda,z,\theta\,|\,\lob,\zob) \;=\;
w_z(z,\zob)\,f_A(\theta,\lob,\zob,\lambda,z)\,\lambda\,
\mathbf{1}_{\{\lambda<\lob\}} ,
\label{eq:rho_prj}
\end{equation}
where $w_z(z,\zob)$ is the parabolic photo-$z$ kernel,
$f_A\in[0,1]$ is the closed-form overlap fraction of the target and
projected-halo \textsc{redMaPPer} disks, the factor $\lambda$ converts
``halo present'' into ``halo contributes $\lambda$ counts'', and the
indicator implements the \textsc{redMaPPer} ranking cut.

The \emph{$\Pop[X]$ operator} is the LOS-averaged integral of an
arbitrary $X(M,z,\theta)$ weighted by the halo distribution and
$\rhoprj$:
\begin{equation}
\begin{aligned}
\Pop[X](\lob,\zob) \;\equiv\;
&\int dz\,\frac{dV}{dz\,d\Omega}(z)
 \int dM\,n(M,z)
 \int d\lambda\,P(\lambda\,|\,M,z) \\
&{}\cdot\;2\pi\int d\theta\,\sin\theta\;
 \rhoprj(\lambda,z,\theta\,|\,\lob,\zob)\,X(M,z,\theta\,|\,\zob) .
\end{aligned}
\label{eq:Pop}
\end{equation}

Three specialisations drive the pipeline:
\begin{equation}
\Pop[1] \;=\; \avg{\dprj_{\rm bkg}}(\lob,\zob),\qquad
I_2 \;\equiv\; \Pop[\,b\,\xinl\,],\qquad
I_1 \;\equiv\; \Pop[\,b\,\xinl\,\sigma(\theta)\,],
\label{eq:Pop_specs}
\end{equation}
where $b=b(M,z)$ is the halo bias and $\xinl(z,\zob,\theta)$ is the
nonlinear matter correlation function evaluated at the 3-D comoving
separation $\Delta\chi(z,\zob,\theta)$ between the target and the
projected halo.  By linearity, $\Pop[\,b\,\xinl(1-\sigma)\,]=I_2-I_1$.
Halo exclusion sets $\xinl\equiv 0$ whenever $\Delta\chi<R_{\rm excl}$.

\subsubsection{Bias pipeline}
\label{sec:pipeline}

The two regime amplitudes of Eq.~\eqref{eq:b_sigmoid} are constructed
in two steps: a per-$\ltr$ closure followed by a
latent-marginalisation against the bin posterior
$P(\ltr\,|\,\Delta\lambda_i,\Delta z_j)$.  With the halo-bias aggregate
$\bbhaloij$ from Eq.~\eqref{eq:bhalo} and the three forward
integrals $\Pop[1]$, $I_1$, $I_2$ in hand, the random-aperture
projection contamination follows from substituting
$\bbsel\!\to\!\bbhaloij$ in the linearised identity,
\begin{equation}
\dprj_{\rm RND}(\lob,\zob) \;=\; \Pop[1](\lob,\zob)
                              + \bbhaloij\,I_2(\lob,\zob) .
\label{eq:dprj_rnd}
\end{equation}
At fixed latent $\ltr$, the LSS-regime closure is the empirical
Buzzard relation
\begin{equation}
\bLSS(\lob,\ltr,\zob) \;=\; \bbhaloij\,
  \bigl[\,1 + 0.13\,\delta^{\rm prj}(\lob,\ltr,\zob)\,\bigr] ,
\qquad
\delta^{\rm prj} \;\equiv\; \frac{\lob-\ltr}{\dprj_{\rm RND}} - 1 ,
\label{eq:b_LSS_latent}
\end{equation}
and the correlated-LOS closure follows from the linear $\bbsel$ identity,
\begin{equation}
\bLOS(\lob,\ltr,\zob) \;=\;
\frac{(\lob-\ltr) - \Pop[1] - \bLSS(\lob,\ltr,\zob)\,I_1}{I_2 - I_1} .
\label{eq:b_LOS_latent}
\end{equation}
Both quantities depend explicitly on $\ltr$, which is unobservable;
they are intermediate objects of the pipeline only.  The physically
meaningful, latent-marginalised regime amplitudes that enter the
mixture model of Eq.~\eqref{eq:b_sigmoid} follow from averaging
Eqs.~\eqref{eq:b_LSS_latent}--\eqref{eq:b_LOS_latent} over the bin
posterior,
\begin{equation}
\begin{aligned}
\bbLOSij &\;=\;
\int d\ltr \; \bLOS(\lob,\ltr,\zob)\,
P(\ltr\,|\,\Delta\lambda_i,\Delta z_j) , \\
\bbLSSij &\;=\;
\int d\ltr \; \bLSS(\lob,\ltr,\zob)\,
P(\ltr\,|\,\Delta\lambda_i,\Delta z_j) ,
\end{aligned}
\label{eq:b_marg_lt}
\end{equation}
with the latent-richness posterior set by the bin,
\begin{equation}
P(\ltr\,|\,\Delta\lambda_i,\Delta z_j) \;\propto\;
\mathcal{S}_{\Delta\lambda_i}(\ltr,\ztr)\,\pi(\ltr,\ztr) ,
\qquad
\pi(\ltr,\ztr) \;=\; \int dM\,n(M,\ztr)\,P(\ltr\,|\,M,\ztr) ,
\label{eq:plt_post}
\end{equation}
and $\mathcal{S}_{\Delta\lambda_i}$ from Eq.~\eqref{eq:Slam_def}.
Inserted into Eq.~\eqref{eq:b_sigmoid}, the two
latent-marginalised regime amplitudes
$\bbLOSij,\bbLSSij$ give $\bbselij(\theta)$, the
final input to the lensing profile of the next section.

% ----------------------------------------------------------------------------
\subsection{Two-halo projected lensing profile}
\label{sec:lensing}

\subsubsection{\texorpdfstring{$\avg{\Sprj(R\,|\,\lob,\zob)}$}{Sprj(R|lob,zob)}}

Combining the selection-affected bias of Eq.~\eqref{eq:b_sigmoid} with
the projected halo distribution of Eq.~\eqref{eq:Pop} yields the
two-halo projected surface density of an optically selected stack,
\begin{multline}
\avg{\Sprj(R\,|\,\lob,\zob)} \;=\;
\int d\theta\,\sin\theta\int dz\,\frac{dV}{dz\,d\Omega}\int dM\,n(M,z) \\
\cdot\;\bigl[\,1 + b(M,z)\,\bbselij(\theta)\,\xinl(z,\zob,\theta)\,\bigr]\,
\Smis(R\,|\,M,z,\theta,\zob) .
\label{eq:Sprj}
\end{multline}
The $1$ inside the bracket is the mean (uncorrelated) contribution and
$b\,\bbsel\,\xinl$ is the correlation-function boost.  $\Smis$ is the
azimuth-averaged miscentered NFW surface density of the contaminating
halo, defined in Eq.~\eqref{eq:Smis} below.

\subsubsection{Lensing excess \texorpdfstring{$\avg{\DSprj(R\,|\,\lob,\zob)}$}{DSprj(R|lob,zob)}}

The two-halo lensing observable returned by the pipeline is the radial
excess of Eq.~\eqref{eq:Sprj}, defined in real space as
\begin{equation}
\avg{\DSprj(R\,|\,\lob,\zob)} \;\equiv\;
\bar\Sigma^{\rm prj}(<\!R\,|\,\lob,\zob)
\;-\;
\avg{\Sprj(R\,|\,\lob,\zob)} ,
\label{eq:DSprj_def}
\end{equation}
with the mean-within-$R$ defined as the standard cylindrical average
\begin{equation}
\bar\Sigma^{\rm prj}(<\!R\,|\,\lob,\zob) \;=\;
\frac{2}{R^{2}}\int_{0}^{R} R'\,\avg{\Sprj(R'\,|\,\lob,\zob)}\,dR' .
\label{eq:Sbar}
\end{equation}
Because the operator $R\mapsto\bar\Sigma(<\!R)-\Sigma(R)$ is linear in
the radial argument only, it commutes with the $(\theta,z,M)$ integrals
of Eq.~\eqref{eq:Sprj}: the only substitution required at the NFW
look-up is
\begin{equation}
\Smis(R\,|\,\cdot)\;\longrightarrow\;
\DSmis(R\,|\,\cdot)\equiv \bar\Smis(<\!R\,|\,\cdot)-\Smis(R\,|\,\cdot) ,
\label{eq:DSmis_swap}
\end{equation}
and the same $(\theta,z,M)$ quadrature carries
$\avg{\Sprj}$ to $\avg{\DSprj}$ unchanged.  No Fourier-space short-cut
is used at any stage of the pipeline.

% ----------------------------------------------------------------------------
\subsection{Miscentering}
\label{sec:miscentering}

The contaminating halo is offset from the radius $R$ at which the
target stack is evaluated, so $\Smis$ is the one-halo surface density
integrated over the azimuthal angle $\varphi$:
\begin{equation}
\Smis(R\,|\,M,z,\theta,\zob) \;=\;
\int_{0}^{2\pi} d\varphi\;\Sigma\!\left(R_h\,|\,M,z\right) ,
\qquad
R_h \;=\; \sqrt{R^{2}+R_\theta^{2}-2 R R_\theta\cos\varphi} ,
\label{eq:Smis}
\end{equation}
with $R_\theta\equiv\theta\,D_A(\zob)$ the physical projected
distance from the target to the contaminating halo, and
$\Sigma(R_h\,|\,M,z)$ the centred NFW surface density of the
contaminating halo at 2-D radius $R_h$ from its own centre.  $\Smis$
is tabulated from the offset-NFW look-up.

In addition, mis-identification of the target's central galaxy
displaces a fraction $f_{\rm mis}$ of the stack itself by a 2-D offset
$R'$ drawn from a Rayleigh distribution,
\begin{equation}
P_{\rm mis}(R'\,|\,\Rmis) \;=\;
\frac{R'}{\Rmis^{2}}\,\exp\!\left[-\frac{R'^{2}}{2\Rmis^{2}}\right] ,
\label{eq:Pmis}
\end{equation}
applied to the one-halo term of Eq.~\eqref{eq:sigma_split} as the
azimuth- and offset-averaged convolution
\begin{equation}
\Sigma^{\rm 1h}_{\rm mis}(R) \;=\;
\int_{0}^{\infty} dR'\,P_{\rm mis}(R'\,|\,\Rmis)
\int_{0}^{2\pi} \frac{d\varphi}{2\pi}\,
\Sigma^{\rm 1h}\!\left(\sqrt{R^{2}+R'^{2}-2RR'\cos\varphi}\right) ,
\label{eq:S1h_mis}
\end{equation}
and the two-population mixture
$\Sigma^{\rm 1h,\,tot}=(1-f_{\rm mis})\,\Sigma^{\rm 1h}+
 f_{\rm mis}\,\Sigma^{\rm 1h}_{\rm mis}$.
The two-halo $\avg{\Sprj}$ inherits target-miscentering through the
same mixture, but in our default pipeline this is absorbed inside the
stacked one-halo term and the two-halo term is computed at the target
truth centre.

% ----------------------------------------------------------------------------
\subsection{\texorpdfstring{Cluster photo-$z$ scatter}{Cluster photo-z scatter}}
\label{sec:photoz}

The photo-$z$ kernel $w_z(z,\zob)$ that enters the projection weight
$\rhoprj$ of Eq.~\eqref{eq:rho_prj} is the parabolic kernel of
Costanzi et~al.\ (2019) calibrated against \textsc{redMaPPer} DR8
photo-$z$ residuals.  We do \emph{not} additionally smooth $\avg{\Sprj}$
by a Gaussian in $\zob$: in the production pipeline target photo-$z$
errors are modelled at the abundance level only (cluster bin migration
in $\zob$), and the impact on $\avg{\DSprj}$ has been verified to be
$<0.4\%$ for typical \textsc{redMaPPer} photo-$z$ uncertainties at
$z\le 0.33$.  The $w_z$ kernel itself sets two of the most important
features of the $\Pop$ operator's $z$-axis structure (the ring-plateau
band and the exclusion peaks of Sec.~\ref{sec:tests:numerics}).

% ----------------------------------------------------------------------------
\subsection{Parameter inference setup}
\label{sec:inference}

The cosmological and nuisance parameters varied in the production MCMC
are listed in Table~\ref{tab:params}.  We use a Gaussian likelihood on
the joint data vector
$\mathbf{d}=\{N(\lob,\zob),\,\avg{\DSprj}(R\,|\,\lob,\zob)\}$ with a
block-diagonal covariance: a Poisson-plus-sample-variance abundance
block, and a shape-noise-dominated lensing block estimated from
jackknife resampling of the survey area.  No global multiplicative
boost on $\avg{\DSprj}$ is included; the projection bias is captured
entirely through $\bbsel(\theta)$ and the resulting two-halo integrand
of Eq.~\eqref{eq:Sprj}.

% ----------------------------------------------------------------------------
\subsection{Calibrated parameters of the production pipeline}
\label{sec:fits}

Table~\ref{tab:params} lists the default parameter values adopted in
the \texttt{mock\_mcmc\_cp\_camb.ini} configuration.  All values are
given for the SDSS \textsc{redMaPPer} redshift range
$0.1\le z\le 0.33$.

\begin{table}[H]
\centering
\small
\begin{tabular}{@{}lll@{}}
\toprule
Parameter & Value & Source \\
\midrule
\multicolumn{3}{@{}l}{\emph{EMG richness kernel, Eq.~\eqref{eq:emg_kernel}}}\\
$\Delta\mu(z\!=\!0.2)$ & calibrated per $\ltr$ bin & Costanzi et~al.\ 2019 \\
$\sigma(z\!=\!0.2)$    & calibrated per $\ltr$ bin & Costanzi et~al.\ 2019 \\
$f^{\rm prj}(z\!=\!0.2)$ & $0.18$ & Costanzi et~al.\ 2019 \\
$\tau(z\!=\!0.2)$        & $0.16$ & Costanzi et~al.\ 2019 \\
\midrule
\multicolumn{3}{@{}l}{\emph{Sigmoid bias, Eq.~\eqref{eq:b_sigmoid}}}\\
$k$        & $2.5/\theta_\lambda$ & Costanzi et~al.\ 2026 \\
$\theta_0$ & $\theta_\lambda/2$   & Costanzi et~al.\ 2026 \\
\midrule
\multicolumn{3}{@{}l}{\emph{Empirical LSS-regime closure, Eq.~\eqref{eq:b_LSS_latent}}}\\
$0.13$ coefficient & Buzzard calibration & Costanzi et~al.\ 2026 \\
\midrule
\multicolumn{3}{@{}l}{\emph{Miscentering, Eq.~\eqref{eq:Pmis}}}\\
$f_{\rm mis}$ & $0.22\pm 0.11$ & Zhang et~al.\ 2019 \\
$\Rmis$       & $(0.29\pm 0.04)\,R_\lambda$ & Zhang et~al.\ 2019 \\
\bottomrule
\end{tabular}
\caption{Default values of the projection-effects parameters used by
the DES Y3 cluster-cosmology pipeline.}
\label{tab:params}
\end{table}

% ============================================================================
\section{Tests and validation}
\label{sec:tests}

\subsection{CAMB vs.\ emulator predictions}
\label{sec:tests:emulator}

Our pipeline builds the nonlinear matter correlation function
$\xinl(r,z)$ from a CAMB linear $P(k,z)$ passed through halofit and an
FFTlog Hankel transform.  We cross-check the resulting $\xinl$ and the
downstream $\avg{\Sprj}$, $\avg{\DSprj}$ against an external emulator
prediction (e.g.\ \texttt{DarkEmulator}, Nishimichi et~al.\ 2019),
using matched cosmologies on a grid in the parameter space probed by
the production MCMC.  At the level of the $\Pop$-operator integrals
$\Pop[1], I_1, I_2$, agreement is required at the
$\sim\!2\%$ level (set by emulator accuracy on $\xi_{\rm hm}$ and
$\xi_{\rm mm}$).  At the level of $\avg{\DSprj}(R\,|\,\lob,\zob)$,
agreement is required at $<\!1\%$ across $0.3\le R\le 30\,\hMpc$ and
$0.1\le\zob\le 0.6$.  The CAMB pipeline is the production default
because it (i) extends naturally outside the emulator training box,
and (ii) does not require interpolation in cosmology at every MCMC
step.

\subsection{Numerical pipeline tests}
\label{sec:tests:numerics}

The three forward integrals $\Pop[1], I_1, I_2$ of
Eq.~\eqref{eq:Pop} are 4-D integrals over $(z,\theta,M,\lambda)$.  The
inner $(M,\lambda)$ axes are smooth and converge under a plain
Gauss--Legendre rule with $N_M=24$, $N_\lambda=60$ to better than
$10^{-4}$.  The outer $(z,\theta)$ axes contain two physics-driven
features that punish any uninformed fixed grid:

\begin{itemize}
\item the \emph{halo-exclusion step in $\theta$}: at each $z$ the
  integrand has a hard discontinuity at
  $\theta=\theta_{\rm excl}(z)$ where $\xinl$ is set to zero by the
  exclusion prescription.  We integrate per-$z$ on the smooth interval
  $[\theta_{\rm excl}(z),\,2\theta_\lambda]$ with no mask;
\item the \emph{exclusion-peak + ring-plateau structure in $z$}: the
  photo-$z$ kernel carves out a window of width $\sim 0.1$ around
  $\zob$, and inside that window $\xinl$ produces twin sharp peaks at
  $z=\zob\pm\delta z_{\rm excl}$ on top of a narrow plateau.  We
  split the $z$-axis into a ring band (GL in $z$) and two outer
  half-planes (GL in $u=\ln|\Delta\chi_\parallel|$).
\end{itemize}

At $N_z=80$, $N_\theta=10$ both $\theta$- and $z$-axis splits reach
sub-$0.01\%$ precision on $\Pop[1], I_1, I_2$ vs.\ a
\texttt{scipy.integrate.quad} reference with matched inner GL grids.
The corresponding regression tests live in
\texttt{tests/test\_integrals.py}.

\subsection{End-to-end validation on Buzzard mocks}
\label{sec:tests:buzzard}

We validate the full pipeline (selection bias $\bbsel(\theta)$ +
two-halo $\avg{\Sprj}$ + lensing excess $\avg{\DSprj}$) against the
Buzzard-v1.1 mock cluster catalogue.  At fixed $(\lob,\zob)$ we
recover the mock $\bbsel(\theta)$ at the few-per-cent level and the
mock $\avg{\DSprj}(R)$ at $<1\%$ on $R\ge 1\,\hMpc$, with the residual
at $R\sim R_{\rm excl}$ tracking the ring-plateau region of the
$z$-integrand discussed above.

% ============================================================================
\section{Application to DES Y3 redMaPPer}
\label{sec:application}

Section~\ref{sec:fits} listed the production parameter set; this
section reports the per-bin predictions of $\avg{\DSprj}(R)$ that feed
the DES Y3 cluster-cosmology likelihood, and contrasts them with the
isotropic two-halo limit ($\bbsel\!\to\!\bbhaloij$) to display the
selection bias as a function of $(\lob,\zob)$.

% ============================================================================
\section{Discussion}
\label{sec:discussion}

The closest analogue to the present model in the literature is the
$\Pi(R)$ anisotropic boost of Park et~al.\ (2023) and Sunayama
et~al.\ (2023), which writes
$\avg{\Sigma}=\Pi(R)\,\avg{\Sigma^{\rm iso}}$ for a piecewise-linear
$\Pi(R)$ calibrated against mocks.  Costanzi et~al.\ (2026) instead
calibrate a multiplicative profile fit
$\mathcal{B}(R)\equiv\avg{\Sprj}/\avg{\Sigma^{\rm iso}}-1$ as a
double power law.  Both formulations factor the projection bias out of
the $(\theta,z,M)$ integrand, at the cost of (i) introducing a global
fitting form whose parameters carry no halo-model interpretation, and
(ii) entangling the small- and large-scale contributions inside a
single multiplicative profile.  Our model instead carries the
projection bias inside the $(\theta,z,M)$ integrand of
Eq.~\eqref{eq:Sprj} via the scale-dependent $\bbsel(\theta)$, with two
amplitudes per observed bin ($\bbLOS,\bbLSS$) that admit a direct
halo-model interpretation as the correlated-LOS and LSS regimes of
the selection-affected cluster bias.  No global multiplicative fitting
function on $\avg{\DSprj}$ is used.

% ============================================================================
\section{Summary}
\label{sec:summary}

The forward model in the \texttt{RichnessSelection} module combines
(i) the Costanzi et~al.\ (2019) EMG richness kernel
$P(\lob\,|\,\ltr,z)$ of Eq.~\eqref{eq:emg_kernel},
(ii) the bias pipeline of Sec.~\ref{sec:pipeline}, which produces the
small- and large-$\theta$ plateaus
$\bbLOSij,\bbLSSij$ from the halo-bias aggregate $\bbhaloij$ and the three
$\Pop$-operator integrals $\Pop[1],I_1,I_2$, and assembles the
selection-affected bias $\bbsel(\theta)$ via the mixture-of-regimes
ansatz of
Eq.~\eqref{eq:b_sigmoid},
(iii) the two-halo projected lensing profile $\avg{\Sprj}$ of
Eq.~\eqref{eq:Sprj}, with its real-space radial excess $\avg{\DSprj}$
of Eqs.~\eqref{eq:DSprj_def}--\eqref{eq:DSmis_swap}, and
(iv) the Rayleigh miscentering kernel of Eq.~\eqref{eq:Pmis}.
All radial integrals are evaluated in real space; no Fourier-space
short-cuts and no global multiplicative fitting functions for the
profile bias are used.

\appendix

\section{Notation summary}
\label{app:notation}

\renewcommand{\arraystretch}{1.2}
\begin{table}[H]
\centering
\small
\begin{tabular}{@{}p{0.22\textwidth} p{0.7\textwidth}@{}}
\toprule
Symbol & Meaning \\
\midrule
$\lob,\ltr$ & observed and intrinsic (true) richness \\
$\zob$ & cluster photometric redshift \\
$M$, $n(M,z)$ & halo mass and Tinker-2008 halo mass function \\
$b(M,z)$ & Tinker-2010 peak-height halo bias \\
$\bbhaloij$ & halo-bias aggregate, bin-averaged, Eq.~\eqref{eq:bhalo} \\
$\bbLOS,\bbLSS$ & correlated-LOS and LSS-regime amplitudes of $\bbsel$ (latent-marginalised), Eq.~\eqref{eq:b_marg_lt} \\
$\bbsel$ & selection-affected cluster bias, Eq.~\eqref{eq:b_sigmoid} \\
$\rhoprj$ & projection kernel weight, Eq.~\eqref{eq:rho_prj} \\
$\Pop[X]$ & LOS-averaged operator, Eq.~\eqref{eq:Pop} \\
$\dprj_{\rm RND}$ & random-aperture projection contamination, Eq.~\eqref{eq:dprj_rnd} \\
$\xinl(z,\zob,\theta)$ & nonlinear matter 2pt correlation, halofit + FFTlog \\
$f^{\rm prj},\tau$ & EMG projection-tail amplitude and inverse scale \\
$\Delta\mu,\sigma$ & EMG background mean bias and width \\
$\avg{\Sprj},\avg{\DSprj}$ & two-halo projected density and lensing excess \\
$\Smis,\DSmis$ & azimuth-averaged miscentered NFW $\Sigma$ and $\Delta\Sigma$ \\
$f_{\rm mis},\Rmis$ & target-miscentering fraction and Rayleigh scale \\
$R_\lambda(\lob)$ & \textsc{redMaPPer} aperture radius, Eq.~\eqref{eq:Rlam} \\
$R_{\rm excl}$ & halo-exclusion radius \\
$\theta_\lambda$ & angle subtended by $R_\lambda$ at $\zob$ \\
\bottomrule
\end{tabular}
\caption{Notation used throughout this document.}
\label{tab:notation}
\end{table}

\end{document}
